% =====================================================================
% CAPÍTULO 1: INTRODUCCIÓN Y PLANTEAMIENTO DEL PROBLEMA
% =====================================================================
\chapter{Introducción y Planteamiento del Problema}
\label{ch:introduccion}

\section{Contexto Institucional}

El \ac{DIC} es la unidad académica responsable de la formación profesional en las
áreas de informática, computación, desarrollo de software y tecnologías de la
información. Su operación diaria involucra cuatro actores fundamentales:

\begin{itemize}[leftmargin=1.4cm, itemsep=4pt]
  \item \textbf{Estudiantes:} demandan cupos, consultan horarios, califican
        docentes, reportan faltas y presentan quejas o sugerencias.
  \item \textbf{Docentes:} administran materias, registran asistencia y faltas,
        se comunican con estudiantes y responden evaluaciones.
  \item \textbf{Secretaría:} recibe, canaliza y da seguimiento a solicitudes,
        quejas y sugerencias; coordina la comunicación entre las partes.
  \item \textbf{Departamento (Dirección/Administración):} verifica resultados,
        analiza información y toma decisiones académicas.
\end{itemize}

A pesar de la naturaleza tecnológica del departamento, sus procesos
administrativos y académicos no cuentan con un sistema integrado que articule
estos actores, lo cual genera reprocesos, demoras y pérdida de información.

\section{Descripción del Problema}

\begin{cajaProblema}[Problema central]
El \ac{DIC} carece de un mecanismo unificado, trazable y analítico para gestionar
la comunicación, las solicitudes y las evaluaciones entre estudiantes, docentes,
secretaría y dirección, lo que provoca demoras en la respuesta, evaluaciones
incompletas, quejas sin seguimiento y decisiones tomadas sin evidencia.
\end{cajaProblema}

A partir del levantamiento inicial de necesidades se identificaron los siguientes
problemas específicos:

\begin{enumerate}[leftmargin=1.4cm, itemsep=5pt]
  \item \textbf{Solicitud de cupos.} El proceso se realiza de forma manual o por
        canales informales (correo, mensajería, presencial), sin control de
        disponibilidad en tiempo real ni criterios de priorización transparentes.

  \item \textbf{Creación y edición de materias.} Las materias nuevas o sus
        modificaciones (contenido, prerrequisitos, créditos, horarios) no siguen
        un flujo formal de solicitud--aprobación--publicación.

  \item \textbf{Comunicación estudiante $\rightarrow$ docente.} No existe un canal
        institucional que garantice que el mensaje llegue, sea leído y tenga
        respuesta dentro de un plazo razonable.

  \item \textbf{Comunicación docente $\rightarrow$ estudiante.} Las notificaciones
        sobre faltas, cambios de horario o novedades académicas dependen de
        medios externos y no quedan registradas.

  \item \textbf{Registro y notificación de faltas.} Las faltas se registran de
        forma tardía y el estudiante se entera cuando el acumulado ya afecta su
        situación académica.

  \item \textbf{Evaluación de docentes por estudiantes.} La evaluación no se
        aplica a tiempo, la participación es baja y los resultados se consolidan
        tarde, cuando ya no permiten retroalimentación oportuna.

  \item \textbf{Quejas y sugerencias.} Se reciben por canales dispersos, sin
        categorización, sin responsable asignado y sin trazabilidad de estado.

  \item \textbf{Efectividad de la secretaría.} La secretaría actúa como cuello de
        botella: recibe las quejas y solicitudes, pero no cuenta con herramientas
        para priorizar, delegar y comunicarse formalmente con el docente.

  \item \textbf{Verificación de resultados.} La secretaría y el departamento no
        disponen de tableros que les permitan verificar resultados, quejas y
        evaluaciones de forma consolidada.

  \item \textbf{Asignación de horarios.} Los horarios se construyen sin conocer
        las preferencias reales de los estudiantes, generando baja asistencia y
        conflictos de disponibilidad.

  \item \textbf{Ausencia de analítica.} No se explotan los datos históricos para
        anticipar saturación de cupos, docentes con bajo desempeño o patrones de
        quejas recurrentes.
\end{enumerate}

\section{Formulación del Problema}

\begin{cajaNota}[Pregunta de investigación principal]
¿Cómo puede un sistema integrado de gestión, comunicación y analítica de datos
mejorar la oportunidad, la trazabilidad y la calidad de las decisiones en el
\ac{DIC}?
\end{cajaNota}

De la pregunta principal se derivan las siguientes preguntas específicas:

\begin{itemize}[leftmargin=1.4cm, itemsep=3pt]
  \item ¿Cuáles son los procesos críticos del \ac{DIC} que presentan mayor
        demora y menor trazabilidad?
  \item ¿Qué arquitectura de información permite unificar la comunicación entre
        estudiantes, docentes y secretaría sin duplicar canales?
  \item ¿Cómo diseñar un mecanismo de evaluación docente que garantice el
        \textbf{100\%} de aplicación dentro del periodo?
  \item ¿Qué variables predicen la preferencia de horarios de los estudiantes y
        cómo pueden usarse para optimizar la programación académica?
  \item ¿Qué modelo analítico permite detectar anomalías (quejas recurrentes,
        cupos saturados, docentes con baja evaluación) y generar soluciones
        rápidas accionables?
  \item ¿Cómo verificar la efectividad de las acciones tomadas por la secretaría
        y el departamento?
\end{itemize}

\section{Justificación}

\subsection{Justificación académica}
Un departamento de informática debe ser el primer ejemplo de digitalización
institucional. Diseñar y documentar este sistema constituye un caso de estudio
aplicado sobre arquitectura de software, modelado \ac{UML} e ingeniería de datos.

\subsection{Justificación operativa}
La automatización de solicitudes de cupos, la edición de materias, el registro
de faltas y la gestión de quejas reduce el trabajo manual de la secretaría,
elimina el retrabajo y libera tiempo para actividades de mayor valor.

\subsection{Justificación estratégica}
La incorporación de analítica de datos y de un algoritmo de soluciones rápidas
convierte la operación cotidiana en información accionable: permite anticipar
problemas, priorizar intervenciones y medir el impacto de las decisiones.

\subsection{Justificación social}
Un canal formal, trazable y confidencial para quejas y sugerencias mejora la
convivencia académica, protege al estudiante y ofrece al docente
retroalimentación oportuna y justa.

\section{Objetivos}

\subsection{Objetivo general}
\begin{cajaObjetivo}[Objetivo general]
Diseñar la documentación metodológica, funcional y de modelado del \ac{SGDIC},
que integre la gestión de cupos, materias, comunicación, faltas, evaluación
docente, quejas y sugerencias, junto con un componente de analítica de datos que
genere soluciones rápidas y verificables para el \ac{DIC}.
\end{cajaObjetivo}

\subsection{Objetivos específicos}
\begin{enumerate}[leftmargin=1.4cm, itemsep=4pt]
  \item Documentar el problema y los procesos actuales del \ac{DIC} mediante
        técnicas de investigación cualitativa y cuantitativa.
  \item Definir la misión, visión, valores y objetivos estratégicos del sistema.
  \item Especificar los requerimientos funcionales y no funcionales por actor.
  \item Modelar el sistema mediante diagramas \ac{UML}: casos de uso, clases,
        secuencia, actividades y estados.
  \item Diseñar el mecanismo de evaluación docente con control de plazos que
        garantice su aplicación a tiempo.
  \item Diseñar el módulo de cuestionarios de materias y horarios, y el modelo de
        preferencias por probabilidad.
  \item Proponer el algoritmo de soluciones rápidas basado en analítica de datos.
  \item Definir el modelo de verificación de resultados para la secretaría y el
        departamento.
  \item Establecer el cronograma, los riesgos y los indicadores de éxito.
\end{enumerate}

\section{Alcance y Delimitaciones}

\subsection{Alcance}
El proyecto cubre el análisis, la metodología de investigación, la especificación
de requerimientos y el modelado completo del sistema, así como la definición del
algoritmo analítico y del modelo de verificación institucional. Incluye:

\begin{itemize}[leftmargin=1.4cm, itemsep=3pt]
  \item Los cuatro actores: estudiante, docente, secretaría y departamento.
  \item Los once procesos problemáticos identificados en la sección anterior.
  \item Los diagramas \ac{UML} de casos de uso, clases, secuencia, actividades y
        estados, embebidos en este documento.
  \item La especificación del algoritmo de soluciones rápidas y las métricas.
\end{itemize}

\subsection{Delimitaciones}
\begin{itemize}[leftmargin=1.4cm, itemsep=3pt]
  \item No se incluye la implementación completa del sistema productivo, sino su
        documentación y modelado.
  \item El análisis se limita al \ac{DIC}; otras unidades académicas podrían
        requerir adaptaciones.
  \item Los datos usados en los ejemplos son sintéticos y anonimizados, en
        cumplimiento de las normas de protección de datos personales (\ac{PII}).
  \item La integración con el sistema central de la universidad se especifica a
        nivel de interfaz (\ac{API} \ac{REST}), sin detallar su implementación.
\end{itemize}

\section{Hipótesis}

\begin{enumerate}[leftmargin=1.4cm, itemsep=4pt]
  \item \textbf{H1.} Un canal unificado de comunicación reduce en al menos
        $60\%$ el tiempo de respuesta a solicitudes y mensajes respecto al
        proceso manual actual.
  \item \textbf{H2.} Un mecanismo de evaluación con recordatorios automáticos,
        plazos visibles e incentivos incrementa la tasa de respuesta docente
        a más del $90\%$ dentro del periodo.
  \item \textbf{H3.} Un modelo de preferencias por probabilidad sobre los
        cuestionarios de horarios permite asignar horarios con mayor nivel de
        aceptación por parte de los estudiantes.
  \item \textbf{H4.} Un algoritmo de priorización basado en analítica reduce el
        tiempo de resolución de quejas críticas y mejora la percepción de
        efectividad de la secretaría.
  \item \textbf{H5.} La trazabilidad completa (estado, responsable, fechas)
        aumenta la confianza institucional y la participación en los canales
        formales.
\end{enumerate}

\section{Metodología Resumida}

La investigación se aborda con un enfoque \textbf{mixto} (cualitativo y
cuantitativo) y un diseño de \textbf{investigación aplicada} complementado con el
ciclo \ac{CRISP-DM} para el componente de analítica de datos. El detalle completo
se presenta en el Capítulo~\ref{ch:metodologia}.

\section{Estructura del Documento}

\begin{table}[H]
\centering
\caption{Estructura de la documentación}
\label{tab:estructura}
\begin{tabularx}{\textwidth}{@{}c l X@{}}
\toprule
\rowcolor{azulClaro}
\textbf{Cap.} & \textbf{Título} & \textbf{Contenido principal} \\
\midrule
1 & Introducción y planteamiento del problema &
    Contexto, problema, preguntas, justificación, objetivos, alcance, hipótesis. \\
2 & Metodología de investigación &
    Enfoque mixto, diseño, población, instrumentos, CRISP-DM, validez y ética. \\
3 & Misión, visión y objetivos &
    Misión, visión, valores, objetivos estratégicos, \ac{KPI} y mapa estratégico. \\
4 & Requerimientos &
    Actores, requerimientos funcionales y no funcionales, reglas de negocio. \\
5 & Modelado \ac{UML} &
    Casos de uso, clases, secuencia, actividades y estados (diagramas embebidos). \\
6 & Algoritmo de soluciones rápidas &
    Analítica de datos, priorización, predicción de horarios, verificación. \\
7 & Planificación &
    Cronograma, recursos, riesgos, indicadores y conclusiones. \\
\bottomrule
\end{tabularx}
\end{table}

\section{Glosario de Términos del Dominio}

\begin{table}[H]
\centering
\caption{Glosario del dominio académico}
\label{tab:glosario}
\begin{tabularx}{\textwidth}{@{}l X@{}}
\toprule
\rowcolor{verdeClaro}
\textbf{Término} & \textbf{Definición} \\
\midrule
Cupo & Plaza disponible de un estudiante en una materia para un periodo académico. \\
Materia & Unidad curricular con código, créditos, prerrequisitos y horario. \\
Falta & Inasistencia de un estudiante a una sesión de clase programada. \\
Queja & Manifestación formal de insatisfacción sobre un servicio o persona. \\
Sugerencia & Propuesta de mejora no vinculada a una inconformidad. \\
Evaluación docente & Instrumento mediante el cual el estudiante valora el desempeño. \\
Cuestionario de horario & Instrumento que captura preferencias y disponibilidad. \\
Solución rápida & Recomendación accionable generada por el motor analítico. \\
Verificación & Contraste de resultados, quejas y evaluaciones por secretaría y dirección. \\
Trazabilidad & Conocer el estado, responsable e historial de un caso. \\
\bottomrule
\end{tabularx}
\end{table}